% ==STYLE-BEGIN==
\documentclass[UTF8,11pt]{ctexart}
\usepackage{amsmath,amssymb,amsthm}
\usepackage[a4paper,margin=2.4cm,top=2.3cm,bottom=2.3cm]{geometry}
\usepackage{fancyhdr}
\usepackage{enumitem}
\usepackage{hyperref}
\hypersetup{hidelinks}

% ---- 版式 ----
\linespread{1.15}
\setlength{\parskip}{0.35em}
\setlength{\parindent}{1.6em}
\setlist{nosep,leftmargin=2em}

% ---- 页眉 ----
\pagestyle{fancy}
\fancyhf{}
\fancyhead[L]{\small\coursename}
\fancyhead[R]{\small\thepage}
\renewcommand{\headrulewidth}{0.4pt}

% ---- 定理环境：中文用正体，标题加粗，按小节共用一条流水号 ----
\newtheoremstyle{cnthm}{0.7em}{0.7em}{\normalfont}{}{\bfseries}{}{0.5em}{}
\theoremstyle{cnthm}
\newtheorem{thm}{定理}[section]
\newtheorem{lem}[thm]{引理}
\newtheorem{prop}[thm]{命题}
\newtheorem{cor}[thm]{推论}
\newtheorem{defn}[thm]{定义}
\newtheorem{rem}[thm]{注}
\newtheorem{exa}[thm]{例}
\newtheorem{conj}[thm]{猜想}
\renewcommand{\proofname}{\bfseries 证明}
\providecommand{\coursename}{}
% ==STYLE-END==
\renewcommand{\coursename}{数学分析\,·\,0916}
\title{数学分析 2024-09-16}
\usepackage{graphicx}
\newtheorem{theorem}{定理}
\newtheorem{example}{例}
\date{}
\begin{document}
\maketitle

\section{整数的算术}

整数对加法和乘法封闭。从加法的角度，可以用反复加 \(1\)
描述正整数；从乘法的角度，大于 \(1\) 的正整数可以分解为素数的乘积。
素数是仅有 \(1\) 和自身两个正因数的正整数。

\section{带余除法}

对 \(n\in\mathbb Z\)、\(m\in\mathbb Z\setminus\{0\}\)，
存在唯一的 \(q,r\in\mathbb Z\)，使
\[
n=qm+r,\qquad 0\leq r<|m|.
\]
其中 \(r\) 是余数。

\section{素数有无穷多个}

反设素数只有 \(p_1,\ldots,p_N\)。
构造
\[
M=p_1p_2\cdots p_N+1.
\]
\(M>1\)，因而有素因数；但对每个 \(p_i\)，
\(M\) 除以 \(p_i\) 的余数均为 \(1\)，矛盾。
故素数有无穷多个。

每个大于等于 \(2\) 的正整数都可分解为素数的乘积。
可用归纳法证明：若 \(N\) 为素数，结论显然；
若 \(N=n_1n_2\) 且 \(1<n_1,n_2<N\)，
对 \(n_1,n_2\) 分别使用归纳假设，再将其素因子合并。

\paragraph{补充：素数定理.}
记 \(\pi(x)\) 为不超过 \(x\) 的素数个数，则
\[
\lim_{x\to+\infty}\frac{\pi(x)}{x/\log x}=1,
\qquad
\lim_{x\to+\infty}\frac{\pi(x)}x=0.
\]
其中 \(\log x\) 为自然对数。

\section{素数平方根的无理性}

设 \(p\) 为素数。反设
\(\sqrt p=n/m\)，其中 \(n,m\) 互素且 \(m\ne0\)。
由
\[
n^2=pm^2
\]
可知 \(p\mid n^2\)，故 \(p\mid n\)。设 \(n=pk\)，
代入得
\[
p^2k^2=pm^2,\qquad m^2=pk^2.
\]
于是 \(p\mid m\)，与 \(n,m\) 互素矛盾。因此
\(\sqrt p\notin\mathbb Q\)；特别地，
\(\sqrt2\notin\mathbb Q\)。

\section{有理数与小数展开}

\[
\mathbb Q
=
\left\{\frac nm:n,m\in\mathbb Z,\ m\ne0\right\}.
\]
有理数的小数展开是有限小数或无限循环小数；
反过来，这两类小数都是有理数。

例如，有限小数可以写成整数与 \(10^k\) 的比。
设循环节所对应的整数为 \(B\)，长度为 \(m\)，
则纯循环小数的循环部分可以表示为
\[
B\sum_{j=1}^{\infty}10^{-jm}
=
\frac{B\,10^{-m}}{1-10^{-m}}\in\mathbb Q.
\]
若前面还有 \(k\) 位不循环的小数，再乘 \(10^{-k}\)
并加上有限小数部分即可。

反过来，对 \(\frac nm\) 作带余除法时，余数只有有限种可能。
余数一旦重复，之后的商位便重复；若出现余数 \(0\)，
则小数展开终止。

\section{补充集合论：映射}

设 \(X,Y\) 为集合。映射 \(f:X\to Y\) 是指：
对每个 \(x\in X\)，都存在唯一的 \(y\in Y\) 与之对应，
记为 \(y=f(x)\)。

\(X\) 称为定义域，\(Y\) 称为靶集，而
\[
f(X)=\{f(x):x\in X\}
\]
称为像集。若 \(y=f(x)\)，则 \(y\) 是 \(x\) 的像，
\(x\) 是 \(y\) 的一个原像。

\begin{itemize}
\item \textbf{单射：}
\(f(x_1)=f(x_2)\Rightarrow x_1=x_2\)。
\item \textbf{满射：}
对每个 \(y\in Y\)，存在 \(x\in X\)，使 \(f(x)=y\)。
\item \textbf{双射：}既是单射又是满射。
\end{itemize}

\end{document}

\end{document}
